\setauthor{AC}
\section{Aktuelle Herausforderungen bei Streaming-Abonnements}
Die Nutzung digitaler Streaming-Dienste hat sich in den letzten Jahren stark verändert. Während zu Beginn oft einzelne Plattformen ausreichten, greifen viele Haushalte heute parallel auf mehrere kostenpflichtige Dienste zu. Gleichzeitig steigen Preise, Inhalte wechseln zwischen Plattformen und die Transparenz über tatsächliche Nutzung und laufende Kosten nimmt ab. Aktuelle Marktbeobachtungen zeigen, dass Haushalte mehrere Streaming-Abonnements parallel verwenden und ihre Ausgaben zunehmend bewusster hinterfragen \cite{deloitte_digital_media_2025}. Damit entsteht ein Problemfeld, das nicht nur mediale Verfügbarkeit, sondern auch organisatorische und finanzielle Aspekte umfasst.

Für Anwender bedeutet dies, dass der Zugang zu Filmen nicht mehr nur von persönlichem Interesse abhängt, sondern zunehmend von der Frage, welche Plattform aktuell abonniert wurde und zu welchem Preis. Je mehr Dienste gleichzeitig genutzt werden, desto schwieriger wird es, Kosten, Nutzen und tatsächliche Verwendung in einem ausgewogenen Verhältnis zu halten. Genau an dieser Stelle wird deutlich, dass Streaming nicht nur ein Unterhaltungsangebot ist, sondern auch ein Verwaltungsproblem im Alltag erzeugen kann.

\section{Probleme durch parallele Abos und doppelte Abbuchungen}
Ein wesentliches Problem paralleler Streaming-Abonnements ist die fehlende Übersicht über laufende Kosten. Werden mehrere Dienste gleichzeitig verwendet, verteilen sich Abrechnungen über unterschiedliche Zeitpunkte, Preisstufen und Vertragsmodelle. Für Benutzer ist dadurch oft nur schwer erkennbar, welche Gesamtkosten pro Monat tatsächlich entstehen und welche Plattformen im Verhältnis dazu regelmäßig genutzt werden. Gerade kleine Einzelbeträge wirken zunächst unproblematisch, summieren sich jedoch über längere Zeiträume zu einem merkbaren Kostenfaktor.

Hinzu kommt, dass unnötige oder doppelte Abonnements häufig nicht unmittelbar auffallen. Benutzer behalten Dienste oft länger aktiv, obwohl bestimmte Plattformen über Wochen oder Monate kaum verwendet werden. Zusätzlich kann es vorkommen, dass ähnliche Inhalte auf mehreren Plattformen verfügbar sind und dadurch mehrere Abos bestehen bleiben, obwohl der praktische Mehrwert begrenzt ist. Ohne zentrale Übersicht entsteht somit leicht eine Situation, in der Kosten weiterlaufen, ohne dass ein bewusster Nutzen damit verbunden ist.

Ein weiterer kritischer Punkt betrifft Kündigungsfristen und Abrechnungszyklen. Werden Verlängerungszeitpunkte nicht aktiv beobachtet, verlängern sich Abonnements automatisch und verursachen weitere Abbuchungen. Für Benutzer bedeutet dies nicht nur finanzielle Mehrkosten, sondern auch zusätzlichen organisatorischen Aufwand. Aus fachlicher Sicht zeigt sich hier deutlich, dass eine Lösung benötigt wird, die nicht nur Anbieter auflistet, sondern auch die Verwaltung persönlicher Abo-Informationen strukturiert unterstützt.

Die daraus entstehenden Problemfelder lassen sich insbesondere in folgende Punkte gliedern:
\begin{itemize}
    \item fehlende Transparenz über monatliche Gesamtkosten,
    \item parallele Zahlungen für wenig oder gar nicht genutzte Dienste,
    \item mangelnde Übersicht über Verlängerungs- und Kündigungszeitpunkte,
    \item fehlender Zusammenhang zwischen aktivem Abo und tatsächlich gesuchten Filmen.
\end{itemize}

\section{Bestehende Lösungen und deren Schwächen}
Für Teilaspekte dieses Problems existieren bereits unterschiedliche Lösungsansätze. Einerseits gibt es allgemeine Finanz- oder Haushaltsanwendungen, mit denen wiederkehrende Zahlungen dokumentiert werden können. Diese Werkzeuge konzentrieren sich jedoch meist auf Kostenkontrolle und berücksichtigen nicht, auf welcher Plattform konkrete Filme verfügbar sind. Andererseits existieren Dienste und Datenbanken, die Filminformationen und Anbieter-Verfügbarkeiten anzeigen. Solche Angebote helfen bei der Suche nach Inhalten, enthalten jedoch in der Regel keine persönliche Verwaltung aktiver Abonnements.

Gerade die Trennung dieser beiden Perspektiven stellt eine wesentliche Schwäche bestehender Lösungen dar. Für Benutzer ist nicht nur relevant, ob ein Film irgendwo verfügbar ist, sondern ob er innerhalb der eigenen bereits bezahlten Abonnements abrufbar ist. Ebenso reicht eine reine Kostenübersicht nicht aus, wenn unklar bleibt, welche Dienste im Hinblick auf konkrete Filminhalte überhaupt einen Mehrwert bieten. Eine Anwendung, die beide Ebenen miteinander verbindet, adressiert daher ein real existierendes praktisches Defizit.

\section{Abgrenzung zu Standard-Webprojekten}
Im Unterschied zu einem klassischen Informations- oder Demo-Webprojekt bestand die Herausforderung dieser Arbeit nicht nur darin, Daten darzustellen oder eine Benutzeroberfläche zu gestalten. Vielmehr musste eine Anwendung konzipiert werden, die externe Filmdaten, Anbieterinformationen, persönliche Abonnements und geschützte Benutzerbereiche sinnvoll zusammenführt. Dadurch entsteht ein deutlich höherer fachlicher Anspruch, weil allgemeine Mediendaten und benutzerspezifische Informationen innerhalb einer gemeinsamen Logik verarbeitet werden müssen.

Die Anwendung \glqq Mediahub\grqq{} grenzt sich insbesondere dadurch ab, dass sie nicht nur eine Filmsuche oder eine einfache Abo-Liste bereitstellt, sondern die zentrale Verwaltung aller Abonnements in den Mittelpunkt stellt. Ergänzend dazu werden Filmdaten mit Verfügbarkeiten verknüpft, sodass Benutzer einen unmittelbaren Bezug zwischen ihren aktiven Diensten und konkreten Inhalten erhalten. Auch die Integration geschützter Bereiche über Keycloak zeigt, dass es sich nicht um eine reine Frontend-Demo handelt, sondern um eine fachlich strukturierte Anwendung mit klarer Trennung zwischen öffentlichen und persönlichen Daten.

Die fachliche Komplexität des Projekts zeigt sich damit vor allem in drei Bereichen:
\begin{itemize}
    \item Verknüpfung externer Filmdaten mit benutzerspezifischen Informationen,
    \item sichere Trennung zwischen öffentlichen und geschützten Funktionen,
    \item praktische Ausrichtung auf einen realen Alltagsnutzen statt auf reine Demonstrationszwecke.
\end{itemize}
